\documentclass[12pt,a4paper]{report}
\usepackage[french]{babel}
\usepackage[margin=2cm]{geometry}
\usepackage{hyperref}
\usepackage{graphicx}
\usepackage{listings}
\usepackage{xcolor}

\lstset{
    basicstyle=\ttfamily\small,
    breaklines=true,
    prebreak=\raisebox{0ex}[0ex][0ex]{\ensuremath{\hookleftarrow}},
    frame=single,
    backgroundcolor=\color{lightgray!20}
}

\title{Application de Gestion de Stock\\Rapport de Stage - Été 2026}
\author{}
\date{\today}

\begin{document}

\maketitle
\newpage
\tableofcontents
\newpage

% ============================================
% CHAPITRE 1: CADRE DU PROJET
% ============================================
\chapter{Cadre du projet}

\section{Introduction}

Ce rapport présente le travail réalisé durant ce stage d'été 2026. Le projet porte sur le développement et la mise en place d'une application complète de gestion de stock incluant une architecture microservices robuste, une infrastructure de monitoring et une pipeline de déploiement continu.

L'objectif principal était de concevoir et déployer une solution complète d'entreprise avec un focus particulier sur la qualité du code, la sécurité et la performance.

\section{1.1 Contexte}

\subsection{Environnement professionnel}

Les entreprises de distribution et de logistique font face à plusieurs défis majeurs dans la gestion de leurs opérations :

\begin{itemize}
    \item Croissance rapide et besoin de scalabilité
    \item Complexité des chaînes d'approvisionnement
    \item Exigences de conformité et de traçabilité
    \item Pression concurrentielle pour l'innovation
\end{itemize}

\subsection{Contexte technologique}

Le contexte technologique actuel favorise :

\begin{itemize}
    \item L'adoption du cloud et de la containerisation
    \item L'automatisation des processus DevOps
    \item La mise en place du monitoring en temps réel
    \item La sécurité intégrée (SecOps)
\end{itemize}

\section{1.2 Problématique}

\subsection{Défis DevOps identifiés}

Les organisations modernes font face à plusieurs défis majeurs dans leur infrastructure informatique :

\begin{itemize}
    \item \textbf{Absence de monitoring} : Pas de visibilité sur la performance et la santé des systèmes
    \item \textbf{Risques de sécurité} : Absence de scanner de vulnérabilités et de contrôle qualité
    \item \textbf{Processus de déploiement non automatisé} : Déploiements manuels, coûteux et sujets à erreurs
    \item \textbf{Manque de traçabilité} : Difficulté de suivre les changements et les déploiements
    \item \textbf{Intégration insuffisante} : Friction entre développement et opérations
    \item \textbf{Absence d'analyse de qualité} : Pas de visibilité sur la qualité du code en production
    \item \textbf{Scalabilité limitée} : Infrastructure non prête pour la croissance
\end{itemize}

\subsection{Questions de recherche}

\begin{enumerate}
    \item Comment concevoir une solution moderne et scalable ?
    \item Comment automatiser le cycle de déploiement ?
    \item Comment intégrer la sécurité dans le processus ?
    \item Comment assurer le monitoring continu ?
\end{enumerate}

\section{1.3 Solution DevOps proposée}

\subsection{Infrastructure DevOps}

Pour résoudre ces défis, une infrastructure DevOps complète a été mise en place avec les composants suivants :

\begin{itemize}
    \item \textbf{Contrôle de version} : Git et GitHub
    \item \textbf{Containerisation} : Docker et Docker Compose
    \item \textbf{CI/CD} : Pipeline Jenkins automatisée
    \item \textbf{Analyse statique} : SonarQube
    \item \textbf{Scan de sécurité} : Trivy
    \item \textbf{Monitoring} : Prometheus et Grafana
    \item \textbf{Alerting} : Règles d'alerte automatiques
    \item \textbf{Orchestration} : Docker Compose (préparation Kubernetes)
\end{itemize}

\subsection{Objectifs DevOps}

La solution DevOps vise à :

\begin{enumerate}
    \item Automatiser complètement la pipeline de déploiement
    \item Garantir la sécurité du code et des images
    \item Maintenir la qualité du code en continu
    \item Monitorer et alerter en temps réel
    \item Réduire les délais de déploiement
    \item Améliorer la traçabilité et la reproductibilité
\end{enumerate}

\newpage

% ============================================
% CHAPITRE 2: CONCEPTS DE BASE (VIRTUALISATION ET CLOUD COMPUTING)
% ============================================
\chapter{Concepts de base (Virtualisation et Cloud Computing)}

\section{Introduction}

Ce chapitre présente les concepts fondamentaux de virtualisation et de cloud computing qui constituent la base de notre infrastructure DevOps.

\section{2.1 La virtualisation}

\subsection{2.1.1 Présentation}

La virtualisation est une technologie permettant de créer des ressources informatiques abstraites plutôt que physiques. Elle permet d'utiliser une même ressource physique pour plusieurs environnements isolés.

\subsubsection{Définition}

La virtualisation est le processus de création d'une version logicielle (virtuelle) d'une ressource physique, comme un serveur, un système d'exploitation ou un périphérique de stockage.

\subsection{2.1.2 Techniques de virtualisation système}

\subsubsection{Hyperviseur de Type 1}

Exécuté directement sur le matériel sans système d'exploitation hôte :

\begin{itemize}
    \item VMware ESXi
    \item Hyper-V (Microsoft)
    \item KVM (Kernel-based Virtual Machine)
\end{itemize}

\subsubsection{Hyperviseur de Type 2}

Exécuté sur un système d'exploitation hôte :

\begin{itemize}
    \item VMware Workstation
    \item VirtualBox
    \item Parallels Desktop
\end{itemize}

\subsection{2.1.3 Types de virtualisation}

\begin{itemize}
    \item \textbf{Virtualisation complète} : Machine virtuelle dispose de tous les éléments matériels
    \item \textbf{Paravirtualisation} : Machine virtuelle consciente de sa nature virtuelle
    \item \textbf{Virtualisation au niveau du système d'exploitation} : Partage du noyau
\end{itemize}

\subsection{2.1.4 La différence entre VM et CT}

\subsubsection{Machines Virtuelles (VM)}

\begin{itemize}
    \item Émulation complète du matériel
    \item Chaque VM a son propre système d'exploitation
    \item Isolation complète
    \item Consommation de ressources plus importante
\end{itemize}

\subsubsection{Conteneurs (CT)}

\begin{itemize}
    \item Partage du noyau du système d'exploitation
    \item Légers et rapides à démarrer
    \item Utilisation efficace des ressources
    \item Isolation au niveau du processus
\end{itemize}

\subsection{Comparaison}

\begin{table}[h]
\centering
\begin{tabular}{|l|l|l|}
\hline
\textbf{Aspect} & \textbf{VM} & \textbf{Conteneur} \\
\hline
Taille & 100 MB - plusieurs GB & 10 MB - 100 MB \\
Démarrage & Secondes à minutes & Millisecondes à secondes \\
Isolation & Complète & Au niveau processus \\
OS & Chacun son OS & Partagé \\
Performance & Surcharge d'hyperviseur & Natif \\
\hline
\end{tabular}
\end{table}

\section{2.2 Cloud Computing}

\subsection{2.2.1 Présentation}

Le cloud computing est un modèle de fourniture de services informatiques via Internet, où les ressources sont distribuées, flexibles et accessibles à la demande.

\subsection{2.2.2 Caractéristiques du Cloud Computing}

Le NIST (National Institute of Standards and Technology) définit 5 caractéristiques essentielles :

\begin{enumerate}
    \item \textbf{Accès à la demande} : Ressources disponibles sans interaction humaine
    \item \textbf{Accès réseau large} : Disponible sur le réseau
    \item \textbf{Mutualisation des ressources} : Ressources partagées entre plusieurs clients
    \item \textbf{Élasticité rapide} : Ressources provisionnées et libérées rapidement
    \item \textbf{Mesure du service} : Consommation mesurée et transparente
\end{enumerate}

\subsection{2.2.3 Les services Cloud}

\subsubsection{IaaS (Infrastructure as a Service)}

\begin{itemize}
    \item Fournisseur : AWS, Azure, Google Cloud, Digital Ocean
    \item Offre : Serveurs virtuels, stockage, réseau
    \item Responsabilité client : OS, Applications, Données
\end{itemize}

\subsubsection{PaaS (Platform as a Service)}

\begin{itemize}
    \item Fournisseur : Heroku, Google App Engine, AWS Elastic Beanstalk
    \item Offre : Runtime, Outils de développement, Bases de données
    \item Responsabilité client : Applications, Données
\end{itemize}

\subsubsection{SaaS (Software as a Service)}

\begin{itemize}
    \item Fournisseur : Microsoft Office 365, Salesforce, Google Workspace
    \item Offre : Applications complètes
    \item Responsabilité client : Données uniquement
\end{itemize}

\subsection{2.2.4 Les modèles de déploiement du Cloud}

\begin{enumerate}
    \item \textbf{Cloud Public} : Infrastructure partagée, accessible au public
    \item \textbf{Cloud Privé} : Infrastructure propriétaire, accès restreint
    \item \textbf{Cloud Hybride} : Combinaison de cloud public et privé
    \item \textbf{Cloud Communautaire} : Infrastructure partagée par un groupe d'organisations
\end{enumerate}

\section{2 Conclusion}

La virtualisation et le cloud computing sont les fondations de notre approche DevOps, permettant la scalabilité, l'flexibilité et l'efficacité des ressources.

\newpage

% ============================================
% CHAPITRE 3: DEVOPS - CONCEPTS ET PRATIQUES ACTUELLES
% ============================================
\chapter{DevOps - Concepts et Pratiques Actuelles}

\section{Introduction}

Ce chapitre explore les concepts fondamentaux de DevOps et leur implémentation à travers les technologies et outils actuels qui transforment la façon dont les organisations développent, déploient et maintiennent leurs applications.

\section{Définition et Philosophie DevOps}

\subsection{Qu'est-ce que le DevOps ?}

DevOps est une approche collaborative qui combine développement logiciel (Dev) et opérations informatiques (Ops). L'objectif est d'automatiser et d'intégrer les processus de développement, de test et de déploiement pour livrer des applications de qualité plus rapidement.

DevOps n'est pas un outil, une technologie ou un rôle. C'est une \textbf{philosophie, une culture et un ensemble de pratiques} qui promeut :

\begin{enumerate}
    \item \textbf{Collaboration} : Communication entre tous les acteurs
    \item \textbf{Automatisation} : Scripts et outils pour réduire l'intervention manuelle
    \item \textbf{Mesure} : Collecte et analyse des métriques
    \item \textbf{Partage} : Transparence et documentation
    \item \textbf{Amélioration continue} : Itération et optimisation
\end{enumerate}

\subsection{Transformation DevOps}

\subsubsection{Avant DevOps}

\begin{itemize}
    \item Séparation claire entre Dev et Ops
    \item Communication limitée entre les équipes
    \item Déploiements lents et manuels
    \item Haute latence entre le code et la production
    \item Responsabilités cloisonnées
\end{itemize}

\subsubsection{Après DevOps}

\begin{itemize}
    \item Collaboration étroite entre Dev et Ops
    \item Automatisation maximale
    \item Déploiements fréquents et rapides
    \item Feedback immédiat
    \item Responsabilité partagée
\end{itemize}

\subsection{Les piliers de DevOps}

\begin{table}[h]
\centering
\begin{tabular}{|l|p{4cm}|}
\hline
\textbf{Pilier} & \textbf{Description} \\
\hline
Planification & Définition des objectifs et des ressources \\
Code & Contrôle de version et collaboration \\
Build & Compilation et empaquetage \\
Test & Vérification de la qualité \\
Release & Préparation du déploiement \\
Deploy & Mise en production \\
Opérer & Maintenance et support \\
Monitor & Supervision et alertes \\
\hline
\end{tabular}
\end{table}

\newpage

\section{Intégration Continue (CI) et Déploiement Continu (CD)}

\subsection{Intégration Continue}

L'intégration continue (CI) est une pratique DevOps où les développeurs intègrent leur code dans un référentiel partagé plusieurs fois par jour. Chaque intégration est vérifiée par un build automatisé et des tests.

\subsubsection{Processus CI}

\begin{enumerate}
    \item \textbf{Commit} : Développeur pousse son code dans Git
    \item \textbf{Webhook Trigger} : Le serveur CI détecte le changement
    \item \textbf{Checkout} : Récupération du code
    \item \textbf{Build} : Compilation et construction
    \item \textbf{Test} : Exécution des tests automatisés
    \item \textbf{Analyse} : SonarQube, Linting, Coverage
    \item \textbf{Report} : Génération de rapports
    \item \textbf{Notification} : Feedback au développeur
\end{enumerate}

\subsubsection{Bénéfices de CI}

\begin{itemize}
    \item \textbf{Qualité} : Détection précoce des bugs
    \item \textbf{Rapidité} : Feedback immédiat
    \item \textbf{Fiabilité} : Processus reproductible
    \item \textbf{Confiance} : Déploiements sûrs
    \item \textbf{Efficacité} : Réduction du temps de cycle
\end{itemize}

\subsection{Déploiement Continu et Distribution Continue}

\subsubsection{Continuous Delivery (CD)}

Capacité à déployer le code en production à tout moment, mais pas nécessairement le faire automatiquement. Représente une préparation complète pour la production.

\subsubsection{Continuous Deployment (CD)}

Déploiement automatique et continu en production après les tests réussis. Le code se déploie sans intervention manuelle.

\subsection{Pipeline CI/CD Complète}

\begin{verbatim}
Code Repository (Git)
  |
  v
Webhook Trigger
  |
  v
Jenkins CI/CD Pipeline
  |
  +---> Code Checkout
  +---> Build Stage
  +---> Unit Tests
  +---> Code Analysis (SonarQube)
  +---> Security Scan (Trivy)
  +---> Docker Build
  +---> Push to Registry
  |
  v
Continuous Delivery / Deployment
  |
  +---> Staging Environment
  +---> Production Environment
  +---> Health Checks
\end{verbatim}

\subsection{Outils de CI/CD - Étude comparative}

\subsubsection{Jenkins}

\begin{itemize}
    \item \textbf{Avantages} : Libre, extensible, largement utilisé, support décentralisé
    \item \textbf{Inconvénients} : Configuration complexe, maintenance requise
    \item \textbf{Utilisation} : Orchestration de pipeline CI/CD
    \item \textbf{Application} : Recommandé pour notre projet
\end{itemize}

\subsubsection{GitLab CI}

\begin{itemize}
    \item \textbf{Avantages} : Intégré, configuration YAML simple, performant
    \item \textbf{Inconvénients} : Moins d'extensions, dépendance GitLab
    \item \textbf{Utilisation} : CI/CD intégré à GitLab
\end{itemize}

\subsubsection{GitHub Actions}

\begin{itemize}
    \item \textbf{Avantages} : Intégré à GitHub, gratuit, facile à configurer
    \item \textbf{Inconvénients} : Limité à GitHub, moins flexible
    \item \textbf{Utilisation} : CI/CD natif GitHub
\end{itemize}

\newpage

\section{Containerisation avec Docker}

\subsection{Concept et Avantages}

Docker est une plateforme de containerisation qui permet de créer, de déployer et d'exécuter des applications dans des conteneurs légers et isolés.

\subsubsection{Pourquoi Docker pour DevOps ?}

\begin{itemize}
    \item \textbf{Portabilité} : Fonctionne partout (dev, test, prod)
    \item \textbf{Légèreté} : Moins de surcharge que les VM
    \item \textbf{Isolement} : Chaque conteneur est complètement isolé
    \item \textbf{Rapidité} : Démarrage en millisecondes
    \item \textbf{Cohérence} : Même environnement partout
    \item \textbf{Reproductibilité} : Infrastructure as Code
\end{itemize}

\subsection{Composants Docker}

\begin{enumerate}
    \item \textbf{Docker Image} : Template statique avec configuration complète
    \item \textbf{Docker Container} : Instance en exécution d'une image
    \item \textbf{Dockerfile} : Fichier de configuration pour construire une image
    \item \textbf{Docker Registry} : Référentiel centralisé d'images (Docker Hub)
    \item \textbf{Docker Compose} : Orchestration multi-conteneurs locale
\end{enumerate}

\subsection{Docker dans la Pipeline CI/CD}

\begin{enumerate}
    \item Build du code
    \item Création d'image Docker
    \item Scan de sécurité de l'image
    \item Push vers registre
    \item Déploiement du conteneur
    \item Monitoring et logs
\end{enumerate}

\section{Contrôle de Version et Collaboration}

\subsection{Git et GitHub}

\subsubsection{Git}

\begin{itemize}
    \item \textbf{Avantages} : Distribué, rapide, gratuit, standard de l'industrie
    \item \textbf{Utilisation} : Contrôle de version décentralisé
\end{itemize}

\subsubsection{GitHub}

\begin{itemize}
    \item \textbf{Avantages} : Interface intuitive, collaboration facile, gratuit pour publics
    \item \textbf{Utilisation} : Hébergement centralisé et gestion de code
    \item \textbf{Intégration DevOps} : Webhooks pour trigger automatique
\end{itemize}

\section{Analyse de Qualité et Sécurité}

\subsection{Analyse Statique du Code - SonarQube}

SonarQube est un outil d'analyse statique qui évalue automatiquement la qualité du code.

\subsubsection{Métriques SonarQube}

\begin{itemize}
    \item \textbf{Bugs} : Violations de logique détectées
    \item \textbf{Vulnérabilités} : Failles de sécurité
    \item \textbf{Code Smells} : Problèmes de maintenabilité
    \item \textbf{Couverture de Tests} : Pourcentage de code testé
    \item \textbf{Duplication} : Code dupliqué
\end{itemize}

\subsubsection{Intégration dans la Pipeline}

SonarQube s'exécute après le build pour évaluer la qualité avant le déploiement.

\subsection{Scan de Sécurité - Trivy}

Trivy est un scanner de vulnérabilités pour images Docker et dépendances.

\subsubsection{Fonctionnalités}

\begin{itemize}
    \item Scan des vulnérabilités du système d'exploitation
    \item Scan des dépendances (npm, pip, Maven, etc.)
    \item Rapport JSON et HTML
    \item Filtrage par sévérité (CRITICAL, HIGH, MEDIUM, LOW)
\end{itemize}

\subsubsection{Intégration DevOps}

\begin{itemize}
    \item Scan automatique après docker build
    \item Arrêt de la pipeline si vulnérabilités critiques
    \item Archivage des rapports pour traçabilité
    \item Alerting immédiat
\end{itemize}

\section{Monitoring et Observabilité}

\subsection{Prometheus - Collecte des Métriques}

Prometheus est un système de monitoring open-source basé sur les métriques.

\subsubsection{Caractéristiques}

\begin{itemize}
    \item Collecte périodique de métriques (time-series)
    \item Requêtes PromQL pour interroger les données
    \item Stockage local optimisé
    \item Intégration facile avec les applications
\end{itemize}

\subsection{Grafana - Visualisation}

Grafana crée des tableaux de bord interactifs pour visualiser les métriques Prometheus.

\subsubsection{Métriques DevOps Essentielles}

\begin{itemize}
    \item Requêtes par seconde (RPS)
    \item Temps de réponse (latence P50, P95, P99)
    \item Utilisation CPU et mémoire
    \item Taux d'erreur HTTP
    \item Disponibilité des services
    \item Couverture de tests
\end{itemize}

\subsection{Alerting et Notifications}

Règles d'alerte configurables pour :

\begin{itemize}
    \item Taux d'erreur élevé
    \item Latence excessive
    \item Ressources insuffisantes
    \item Services indisponibles
\end{itemize}

\section{Outils d'Infrastructure et Orchestration}

\subsection{Docker Compose}

Orchestration multi-conteneurs pour l'environnement local et léger.

\subsection{Kubernetes (Perspective Future)}

Orchestration avancée pour production scalable et haute disponibilité.

\subsection{Terraform et Infrastructure as Code}

Gestion déclarative de l'infrastructure cloud multi-provider.

\subsection{Ansible}

\begin{itemize}
    \item \textbf{Avantages} : Sans agent, YAML simple, agentless
    \item \textbf{Utilisation} : Configuration et déploiement
\end{itemize}

\newpage

% ============================================
% CHAPITRE 4: ETUDE COMPARATIVE DES OUTILS DEVOPS
% ============================================
\chapter{Étude Comparative et Sélection des Outils}

\section{Introduction}

Ce chapitre synthétise l'étude comparative des outils DevOps sélectionnés pour notre infrastructure, en détaillant les critères d'évaluation et les justifications des choix effectués.

\section{Méthodologie d'évaluation}

\subsection{Critères d'évaluation}

Les outils ont été évalués selon les critères suivants :

\begin{itemize}
    \item \textbf{Fonctionnalités} : Réponse aux besoins du projet
    \item \textbf{Facilité d'utilisation} : Courbe d'apprentissage et ergonomie
    \item \textbf{Performance} : Capacité à traiter les charges
    \item \textbf{Coût} : Licences et infrastructure requise
    \item \textbf{Support communautaire} : Documentation et ressources disponibles
    \item \textbf{Intégration} : Compatibilité avec d'autres outils
    \item \textbf{Scalabilité} : Croissance future
\end{itemize}

\section{Récapitulatif des Outils Sélectionnés}

\subsection{Couche VCS (Version Control System)}

\subsubsection{Solution Retenue : Git + GitHub}

\begin{itemize}
    \item \textbf{Git} : Système de contrôle de version distribué
    \item \textbf{GitHub} : Plateforme hébergée avec webhooks pour CI/CD
\end{itemize}

\subsubsection{Justification}

\begin{itemize}
    \item Standard industrie avec écosystème mature
    \item Intégration native avec Jenkins
    \item Collaboration facilitée via GitHub
    \item Gratuit pour projets publics
\end{itemize}

\subsection{Couche CI/CD}

\subsubsection{Solution Retenue : Jenkins}

\begin{itemize}
    \item \textbf{Avantages} :
    \begin{itemize}
        \item Complètement open-source et gratuit
        \item Très extensible via plugins
        \item Large communauté et documentation
        \item Support décentralisé
        \item Exécution distribuée possible
    \end{itemize}
    \item \textbf{Inconvénients} :
    \begin{itemize}
        \item Configuration complexe nécessitant expertise
        \item Maintenance et gestion de l'instance requises
    \end{itemize}
\end{itemize}

\subsubsection{Alternatives Étudiées}

\begin{itemize}
    \item \textbf{GitLab CI} : Intégré mais plus limité
    \item \textbf{GitHub Actions} : Limité à GitHub uniquement
\end{itemize}

\subsection{Couche Containerisation}

\subsubsection{Solution Retenue : Docker}

\begin{itemize}
    \item Standard de containerisation en 2026
    \item Écosystème complet (Docker Compose, Docker Hub)
    \item Intégration facile avec Jenkins
    \item Portabilité garantie
\end{itemize}

\subsubsection{Orchestration}

\begin{itemize}
    \item \textbf{Court terme} : Docker Compose (léger, suffisant pour stage)
    \item \textbf{Long terme} : Kubernetes (scalabilité entreprise)
\end{itemize}

\subsection{Couche Qualité et Sécurité}

\subsubsection{Analyse Statique : SonarQube}

\begin{itemize}
    \item Leader du marché pour analyse de code
    \item Support multilangage
    \item Intégration pipeline facile
    \item Version gratuite suffisante
\end{itemize}

\subsubsection{Scan de Sécurité : Trivy}

\begin{itemize}
    \item Spécialisé dans les vulnérabilités d'images Docker
    \item Très rapide et léger
    \item Open-source et maintenu activement
    \item Integration simple dans pipeline Docker
\end{itemize}

\subsection{Couche Monitoring}

\subsubsection{Prometheus + Grafana}

\begin{itemize}
    \item \textbf{Prometheus} : 
    \begin{itemize}
        \item Collecte de métriques time-series
        \item Push/pull metrics
        \item Query language puissant (PromQL)
    \end{itemize}
    \item \textbf{Grafana} :
    \begin{itemize}
        \item Visualisation professionnelle
        \item Tableaux de bord customisables
        \item Support d'alerting
    \end{itemize}
\end{itemize}

\subsubsection{Justification du Choix}

\begin{itemize}
    \item Stack open-source reconnue
    \item Performance excellente
    \item Faible overhead ressources
    \item Large adoption industrie
\end{itemize}

\section{Matrice de Comparaison}

\begin{table}[h]
\centering
\begin{tabular}{|l|c|c|c|c|}
\hline
\textbf{Critère} & \textbf{Jenkins} & \textbf{GitLab CI} & \textbf{GitHub Actions} \\
\hline
Coût & Gratuit & Freemium & Gratuit \\
Facilité & Moyen & Facile & Très facile \\
Flexibilité & Très haute & Haute & Moyenne \\
Extensibilité & Très haute & Haute & Moyenne \\
\hline
\end{tabular}
\end{table}

\section{Conclusion du Chapitre}

Les outils sélectionnés forment un stack DevOps cohérent et moderne qui répond aux besoins du projet tout en permettant l'évolution future vers une infrastructure Kubernetes.

\section{2. Outils de gestion de code source}

\subsection{Git}

\begin{itemize}
    \item \textbf{Avantages} : Distribué, rapide, gratuit, standard de l'industrie
    \item \textbf{Inconvénients} : Courbe d'apprentissage
    \item \textbf{Utilisation} : Contrôle de version
\end{itemize}

\subsection{GitHub}

\begin{itemize}
    \item \textbf{Avantages} : Interface intuitive, collaboration facile, gratuit pour projets publics
    \item \textbf{Inconvénients} : Dépendance cloud
    \item \textbf{Utilisation} : Hébergement et gestion de code
\end{itemize}

\section{3. Outils d'intégration continue}

\subsection{Jenkins}

\begin{itemize}
    \item \textbf{Avantages} : Libre, extensible, largement utilisé
    \item \textbf{Inconvénients} : Configuration complexe
    \item \textbf{Utilisation} : Orchestration de pipeline CI/CD
\end{itemize}

\subsection{GitLab CI}

\begin{itemize}
    \item \textbf{Avantages} : Intégré, configuration simple, performant
    \item \textbf{Inconvénients} : Moins d'extensions
    \item \textbf{Utilisation} : CI/CD intégré à GitLab
\end{itemize}

\subsection{GitHub Actions}

\begin{itemize}
    \item \textbf{Avantages} : Intégré à GitHub, gratuit, facile
    \item \textbf{Inconvénients} : Limité à GitHub
    \item \textbf{Utilisation} : CI/CD natif GitHub
\end{itemize}

\section{4. Outils de gestion de configuration et automation}

\subsection{Ansible}

\begin{itemize}
    \item \textbf{Avantages} : Sans agent, YAML simple, agentless
    \item \textbf{Inconvénients} : Plus lent pour grandes infrastrures
    \item \textbf{Utilisation} : Configuration et deployment
\end{itemize}

\subsection{Terraform}

\begin{itemize}
    \item \textbf{Avantages} : Infrastructure as Code, multi-cloud
    \item \textbf{Inconvénients} : Courbe d'apprentissage HCL
    \item \textbf{Utilisation} : Provisionnement d'infrastructure
\end{itemize}

\subsection{Kubernetes}

\begin{itemize}
    \item \textbf{Avantages} : Orchestration puissante, scalabilité
    \item \textbf{Inconvénients} : Complexité élevée
    \item \textbf{Utilisation} : Orchestration conteneurs
\end{itemize}

\newpage

% ============================================
% CHAPITRE 5: CONCEPTION ET RÉALISATION
% ============================================
\chapter{Conception et Réalisation}

\section{Introduction}

Ce chapitre détaille la conception et l'implémentation de notre solution DevOps complète pour l'application de gestion de stock.

\section{2. Objectifs du projet DevOps}

\subsection{Objectifs primaires}

\begin{enumerate}
    \item Automatiser complètement la pipeline CI/CD de bout en bout
    \item Containeriser toutes les applications pour la portabilité
    \item Implémenter une infrastructure de monitoring en temps réel
    \item Intégrer les scans de sécurité dans le processus de déploiement
    \item Analyser la qualité du code automatiquement
    \item Orchestrer les déploiements avec Docker
\end{enumerate}

\subsection{Objectifs secondaires}

\begin{enumerate}
    \item Réduire le temps entre la détection d'une vulnérabilité et sa correction
    \item Fournir une traçabilité complète des déploiements
    \item Créer une infrastructure reproductible et portable
    \item Documenter les processus et les meilleures pratiques
    \item Préparer l'infrastructure pour la scalabilité avec Kubernetes
\end{enumerate}

\section{3. Architecture DevOps globale}

\subsection{Flux de déploiement}

\begin{verbatim}
GitHub Repository
  |
  +--> Webhook trigger
  |
  v
Jenkins Pipeline
  |
  +---> Clone Repository
  +---> SonarQube Analysis
  +---> Build Applications
  +---> Docker Build Images
  +---> Trivy Security Scan
  +---> Push to Registry
  +---> Deploy to Production
  +---> Health Check
  |
  v
Production Environment
  |
  +---> Services (Docker Containers)
  +---> Database (PostgreSQL)
  +---> Monitoring Stack
\end{verbatim}

\subsection{Stack technologique DevOps}

\begin{table}[h]
\centering
\begin{tabular}{|l|l|l|}
\hline
\textbf{Catégorie} & \textbf{Outil} & \textbf{Fonction} \\
\hline
VCS & GitHub & Contrôle de version \\
CI/CD & Jenkins & Orchestration pipeline \\
Analyse statique & SonarQube & Qualité du code \\
Scan sécurité & Trivy & Vulnérabilités images \\
Contenarisation & Docker & Empaquetage applications \\
Orchestration & Docker Compose & Orchestration multi-conteneurs \\
Monitoring & Prometheus & Collecte métriques \\
Visualisation & Grafana & Tableaux de bord \\
\hline
\end{tabular}
\end{table}

\section{Création d'une pipeline CI/CD à l'aide de Jenkins et de Docker}

\subsection{Architecture Jenkins}

\begin{verbatim}
Jenkinsfile
  |
  +---> Clone Repository
  +---> SonarQube Analysis
  +---> Build Applications
  +---> Docker Build
  +---> Security Scan (Trivy)
  +---> Push Registry
  +---> Deploy Production
  +---> Health Check
\end{verbatim}

\subsection{Stages détaillés}

\subsubsection{Stage 1 : Clone Repository}

Récupère le dernier code depuis GitHub.

\subsubsection{Stage 2 : SonarQube Analysis}

Analyse la qualité du code backend et frontend.

\subsubsection{Stage 3 : Build Applications}

Compile les applications et génère les artifacts.

\subsubsection{Stage 4 : Docker Build}

Crée les images Docker pour backend et frontend.

\subsubsection{Stage 5 : Security Scan - Trivy}

Détecte les vulnérabilités dans les images Docker.

\subsubsection{Stage 6 : Push to Registry}

Pousse les images validées vers Docker Hub.

\subsubsection{Stage 7 : Deploy to Production}

Récupère les images et redémarre les services.

\subsubsection{Stage 8 : Health Check}

Vérifie la disponibilité de tous les services.

\section{Configuration Docker Compose pour DevOps}

\begin{lstlisting}
version: '3.8'

services:
  prometheus:
    image: prom/prometheus
    ports:
      - "9090:9090"
    volumes:
      - ./monitoring/prometheus.yml:/etc/prometheus/prometheus.yml
      - prometheus_data:/prometheus
    command:
      - '--config.file=/etc/prometheus/prometheus.yml'
      - '--storage.tsdb.path=/prometheus'

  grafana:
    image: grafana/grafana
    ports:
      - "3000:3000"
    environment:
      GF_SECURITY_ADMIN_PASSWORD: admin
    volumes:
      - grafana_data:/var/lib/grafana
    depends_on:
      - prometheus

  jenkins:
    image: jenkins/jenkins:lts
    ports:
      - "8080:8080"
      - "50000:50000"
    volumes:
      - jenkins_home:/var/jenkins_home
      - /var/run/docker.sock:/var/run/docker.sock
    environment:
      JAVA_OPTS: -Djenkins.install.runSetupWizard=false

volumes:
  prometheus_data:
  grafana_data:
  jenkins_home:
\end{lstlisting}

\section{Monitoring avec Prometheus et Grafana}

\subsection{Configuration Prometheus}

\begin{lstlisting}
global:
  scrape_interval: 15s
  evaluation_interval: 15s

scrape_configs:
  - job_name: 'backend'
    static_configs:
      - targets: ['backend:3001']

  - job_name: 'prometheus'
    static_configs:
      - targets: ['localhost:9090']
\end{lstlisting}

\subsection{Dashboards Grafana}

Métriques affichées en temps réel :

\begin{itemize}
    \item Requêtes par seconde (RPS)
    \item Temps de réponse (latence P50, P95, P99)
    \item Utilisation CPU et mémoire
    \item Taux d'erreur HTTP
    \item Connections base de données
    \item Disponibilité des services
\end{itemize}

\section{Mettre en place l'outil Ansible}

\subsection{Introduction à Ansible}

Ansible est un outil d'automatisation qui permet de gérer plusieurs serveurs sans installer d'agent.

\subsection{Playbooks Ansible}

\begin{lstlisting}
---
- hosts: all
  tasks:
    - name: Install Docker
      apt:
        name: docker.io
        state: present

    - name: Start Docker
      service:
        name: docker
        state: started

    - name: Deploy application
      shell: docker-compose up -d
\end{lstlisting}

\section{Installation et configuration}

\subsection{Prérequis}

\begin{itemize}
    \item Docker Desktop installé
    \item Git configuré
    \item Node.js 20+ installé
    \item PostgreSQL 14+ disponible
\end{itemize}

\subsection{Étapes de déploiement}

\begin{enumerate}
    \item Cloner le repository
    \item Configurer les variables d'environnement
    \item Builder les images Docker
    \item Lancer Docker Compose
    \item Accéder aux applications
\end{enumerate}

\subsection{URLs d'accès aux outils DevOps}

\begin{table}[h]
\centering
\begin{tabular}{|l|l|l|}
\hline
\textbf{Service} & \textbf{URL} & \textbf{Fonction} \\
\hline
Jenkins & http://localhost:8080 & Orchestration CI/CD \\
Prometheus & http://localhost:9090 & Collecte métriques \\
Grafana & http://localhost:3000 & Visualisation monitoring \\
SonarQube & http://localhost:9000 & Analyse code \\
\hline
\end{tabular}
\end{table}

\section{Conclusion}

La mise en œuvre de cette architecture DevOps complète démontre comment intégrer les meilleures pratiques modernes pour créer une pipeline automatisée, sécurisée et performante.

\newpage

% ============================================
% CONCLUSION GÉNÉRALE
% ============================================
\chapter*{Conclusion générale}
\addcontentsline{toc}{chapter}{Conclusion générale}

\section*{Réalisations principales}

Ce stage a permis de mener à bien le développement complet d'une application de gestion de stock moderne avec une infrastructure DevOps professionnelle :

\begin{itemize}
    \item Pipeline CI/CD complètement automatisée avec Jenkins
    \item Containerisation Docker de tous les services
    \item Scan de sécurité automatisé (Trivy)
    \item Analyse statique du code (SonarQube)
    \item Infrastructure de monitoring temps réel (Prometheus + Grafana)
    \item Alertes automatiques pour anomalies
    \item Infrastructure entièrement reproduible
\end{itemize}

\section*{Compétences DevOps développées}

\begin{itemize}
    \item Orchestration Jenkins et pipeline CI/CD
    \item Containerisation Docker et Docker Compose
    \item Sécurité : Scan Trivy et hardening images
    \item Qualité : SonarQube et métriques
    \item Monitoring : Prometheus et Grafana
    \item Infrastructure as Code
    \item Best practices DevOps
\end{itemize}

\section*{Perspectives futures}

L'infrastructure peut être améliorée avec :

\begin{enumerate}
    \item Déploiement Kubernetes
    \item GitOps avec ArgoCD
    \item Infrastructure as Code (Terraform)
    \item Stack de logging ELK
    \item Tests de performance
    \item Blue-Green Deployment
\end{enumerate}

\section*{Apprentissages clés}

\begin{itemize}
    \item L'automatisation économise du temps et réduit les erreurs
    \item La sécurité doit être intégrée dès le départ
    \item Le monitoring continu est crucial pour la stabilité
    \item La traçabilité des déploiements est essentielle
    \item La documentation est vitale pour la continuité
\end{itemize}

\end{document}
